\begin{lemma}[Coprime relative extension without an intersection assumption]
\label{lem:coprime-relative-general}
Let \(G=RH\), with \(R\) closed normal and \(H\) closed.
Let \(E\triangleleft G\) be closed and contained in \(R\), and
suppose every finite quotient of \(R/E\) has order prime to \(r\).
For each finite abelian \(r\)-group \(A\), restriction gives a
bijection between continuous characters \(G\to A\) killing \(H\)
and continuous \(G\)-invariant characters \(E\to A\) killing
\(E\cap H\).
\end{lemma}

\begin{proof}
Let \(\chi:E\to A\) be such a character.
The quotient \(L=R/\ker\chi\) has the central finite
\(r\)-subgroup \(Z=E/\ker\chi\cong\chi(E)\), and \(L/Z\)
is pro-\(r'\).
It has a unique pro-\(r'\) complement to \(Z\).

Here is a direct proof of this splitting assertion.
In a finite central extension of a group \(Q\), of order
\(m\) prime to \(r\), by \(Z\), choose a section with normalized
two-cocycle \(c\). Define
\[
 b(q)=m^{-1}\sum_{t\in Q}c(q,t).
\]
The inverse of \(m\) acts on the finite \(r\)-group \(Z\).
Summing the cocycle identity over its last variable gives
\(c(q,t)=b(q)+b(t)-b(qt)\).
Correcting the section by \(-b\) therefore gives a homomorphic
section. Two such sections differ by a homomorphism \(Q\to Z\),
which is zero. The complement is unique.
For \(L\), pass through the cofinal open normal subgroups
\(U\) with \(U\cap Z=1\).
The unique complements in \(L/U\) are compatible and map onto
each other. Their inverse limit gives \(L=Z\times C\).

The \(G\)-action fixes \(Z\) pointwise and preserves its unique
complement. Projection therefore gives a continuous
\(G\)-invariant extension \(\theta:R\to A\).
It is unique, because \(R/E\) has no nonzero character in \(A\).
The restriction of \(\theta\) to \(R\cap H\) factors through
\((R\cap H)/(E\cap H)\), a closed subgroup of \(R/E\), and
hence vanishes.
Now \(\Phi(rh)=\theta(r)\) is well-defined, is a homomorphism
by normality and invariance, and is continuous by compactness.
It kills \(H\), and uniqueness on \(R\) proves uniqueness on
\(G\). The reverse restriction is immediate.
\end{proof}

\begin{theorem}[A prime-specific residual group]
\label{thm:primary-residual-gate}
For a prime \(r\), put
\[
 E_r=R_D\cap\bigcap_{p\ne r}\ker\Phi_p,
 \qquad X_r=H^1(N,\mathbf F_r)^G,
\]
where \(\Phi_p\) is the action on the free pro-\(p\) group.
Then \(E_r\) is closed and normal in \(G\), and
\begin{equation}
 B_{\max r}\cong
 \left(E_r/\overline{[G,E_r](E_r\cap H)}\right)_{\max r}. \label{eq:primary-residual-1}
\end{equation}
For odd \(r\), restriction
\[
 X_r\longrightarrow H^1(E_r,\mathbf F_r)
\]
is injective.
For \(r=2\), its kernel is
\begin{equation}
 \bigoplus_{q\ {\rm odd}}\mathbf F_2\psi_q
                         \oplus\mathbf F_2\eta_3.       \label{eq:primary-residual-2}
\end{equation}
Moreover \(E_r\) has infinite index in \(R_D\).
\end{theorem}

\begin{proof}
The normalized image \(\Phi_p(M)\) is pro-\(p\).
On each finite algebra
\(\mathbf F_p\langle S,T\rangle/(S,T)^d\), its action is
upper unitriangular because it fixes the associated graded.
These finite \(p\)-group actions separate normalized parameters.
Consequently
\[
 R_D/E_r\lhook\joinrel\longrightarrow
                    \prod_{p\ne r}\Phi_p(R_D)
\]
is pro-\(r'\).
Indeed every finite quotient of the closed image factors
through a finite collection of finite \(p\)-group coordinates,
with all \(p\ne r\).
Normality and closedness of \(E_r\) follow from its definition.

Apply Lemma~\ref{lem:coprime-relative-general} with \(R=R_D\),
using the established factorization \(G=R_DH\).
It identifies all finite abelian \(r\)-valued characters of
the two sides of \eqref{eq:primary-residual-1}.
More directly, the natural map from the right side into
\(B_{\max r}\) is injective because these finite characters
separate points and all extend uniquely.
Its cokernel is a pro-\(r\) quotient of \(R_D/E_r\), so is
trivial. Compactness proves \eqref{eq:primary-residual-1}.

A character of \(R_D\) with values in \(\mathbf F_r\) and
trivial on \(E_r\) factors through the pro-\(r'\) quotient
and is zero. Thus
\begin{equation}
 \ker\bigl(X_r\to H^1(E_r,\mathbf F_r)\bigr)
 =
 \ker\bigl(X_r\to H^1(R_D,\mathbf F_r)\bigr).             \label{eq:primary-residual-3}
\end{equation}
For odd \(r\), a character in the latter kernel first vanishes
on \(D\), since \(D/R_D\) is pro-two.
It factors through the rational Kummer quotient \(N/D\).
Complex conjugation acts on this quotient by inversion.
Invariance therefore gives \(\chi=-\chi\), and hence
\(\chi=0\) in odd characteristic.
For \(r=2\), \eqref{eq:primary-residual-3} gives \eqref{eq:primary-residual-2} from the established original
restriction-kernel theorem.

It remains to prove infinite index.
Choose an odd prime \(p\ne r\).
Only one nonzero arithmetic degree-three class at \(p\) is
needed. Its existence follows from
\cite[Remark~8.6, p.~151]{HainMatsumoto2003}; the outer-action
filtration is defined there in \S1, equation~\eqref{eq:primary-residual-2}, p.~120.
Choose an arithmetic representative \(t\in I_p^3H\) whose
degree-three outer class \(v\) is nonzero over \(\mathbf Q_p\).
In particular \(\lambda_p(t)=1\).

Choose \(h\in H\) with \(u=\lambda_p(h)=1+p\), let
\(b=[h,t]\), and let \(\varepsilon_p\in\widehat{\mathbf Z}\)
have \(p\)-component one and all other components zero.
Set
\begin{equation}
 z=b^{\varepsilon_p},\qquad
 w_p=hzh^{-1}z^{-\lambda(h)}.                            \label{eq:primary-residual-4}
\end{equation}
The commutator \(b\) has full multiplier one.
Thus, for every \(\ell\ne p\), its normalized pro-\(\ell\)
image is killed by the power \(\varepsilon_p\).
Both \(z\) and \(w_p\) have trivial \(\ell\)-action for
\(\ell\ne p\).
All rational Kummer coordinates cancel exactly in \eqref{eq:primary-residual-4},
because their conjugation weight on \(M\) is the full
multiplier \(\lambda(h)\). Therefore \(w_p\in E_p\cap H\).

Conjugation by multiplier \(u\) acts on the degree-three
outer class by \(u^3\); this is also immediate from the
linear scaling of a leading derivation.
The class of \eqref{eq:primary-residual-4} is consequently
\begin{equation}
 \operatorname{gr}_{3,p}(w_p)
      =(u^3-u)(u^3-1)v\ne0.                             \label{eq:primary-residual-5}
\end{equation}
Both factors are nonzero, also when \(p=3\).
Working with the outer class suffices to detect a
nontrivial action and avoids any additional signed
parameter comparison.
Its procyclic image has an infinite \(p\)-adic graded
coordinate.
Since \(E_r\subset\ker\Phi_p\), this gives an infinite
image in \(R_D/E_r\), proving infinite index.
\end{proof}

\begin{remark}
The theorem does not assert \(G=E_rH\).
For odd \(r\), it proves injectivity of the indicated
restriction map; it does not compute \(X_r\).
The strictness argument uses one degree-three arithmetic
class and requires no full leading-lattice result at an
odd prime.
\end{remark}
